\section{Setup}
\label{sec:setup}

\paragraph{Datasets and protocol.}
We use four public datasets (Table~\ref{tab:datasets}): German Credit~\citep{statlog_german} and
Breast Cancer Wisconsin Diagnostic~\citep{wolberg1995wdbc}, two small binary tabular tasks, SST-2
sentiment~\citep{socher2013sst,wang2019glue} and four-class AG News~\citep{zhang2015agnews}. The
text datasets are evaluated on their complete official held-out split, with training and
validation data taken from the official training split only. The tabular datasets are too small
for one split to be reliable, so they use \NRepeats{} seeded, stratified
\SplitFractions{} train/validation/test repeats. All identifiers were frozen, with a content
hash, before any template check, training or threshold selection.

The two kinds of dataset pose different problems. The text test sets are balanced
(\NTestNegativeSst{} negative and \NTestPositiveSst{} positive SST-2 sentences;
\NTestPerClassMinAgnews{} AG News articles in each class) and we treat their errors as equally
costly, so accuracy is a fair summary. The tabular datasets are small and imbalanced
(\PctTestPositiveGerman\% of applicants are bad, \PctTestPositiveWdbc\% of masses malignant),
and their errors are not equally costly: missing a bad applicant or a malignant mass is worse
than a false alarm. For them we therefore report the counts of each kind of error and a decision
cost in which a miss costs five times a false alarm (canonical for German Credit, assumed for
Breast Cancer Wisconsin). Held-out datasets get 95\%
bootstrap intervals; tabular datasets get the mean and standard deviation over the repeats
(Appendix~\ref{app:protocol}).

\begin{table}[t]
  \caption{Datasets. Evaluation examples are the official held-out split (text) or the test part
    of each of \NRepeats{} split repeats (tabular); training and validation sizes are per repeat.}
  \label{tab:datasets}
  \centering\small
  \begin{tabular}{llrrrr}
\toprule
Dataset & Type & Classes & Evaluation & Training & Validation \\
\midrule
German Credit & tabular & 2 & 10 $\times$ 200 & 600 & 200 \\
Breast Cancer Wisconsin & tabular & 2 & 10 $\times$ 114 & 341 & 114 \\
SST-2 & text & 2 & 872 & 65{,}349 & 2{,}000 \\
AG News & text & 4 & 7{,}600 & 115{,}000 & 5{,}000 \\
\bottomrule
\end{tabular}

\end{table}

\paragraph{Jev, zero-shot.}
Jev (version \texttt{\JevModel})~\citep{typesafe_sdk} is queried with one question per example: a
one-sentence instruction and one description per class, in a fixed order. Tabular rows are
serialized as one \texttt{name: value} line per feature; text is sent as is. Each dataset has one
template, frozen without tuning. Appendix~\ref{app:payloads} shows complete requests and
responses.

\paragraph{Qwen3.6-27B, zero-shot.}
The open-weight \texttt{\QwenModelId}~\citep{qwen36} is self-hosted and receives exactly Jev's
content: its instruction, options and descriptions in the same order, then the serialized input.
To obtain a class distribution rather than a generated string, every option is scored as the
answer: the log-probabilities of the tokens of the option name and of the end-of-turn token are
summed, and a softmax over the options gives $\mu_x$ (Appendix~\ref{app:qwen}). Thinking is
disabled, except on \ThinkSubsampleN{} items per dataset reported separately.

\paragraph{XGBoost and BERT, supervised.}
XGBoost~\citep{chen2016xgboost} uses one-hot tabular features or TF-IDF text features; a
\texttt{bert-base} encoder~\citep{devlin2019bert} with a classification head is fine-tuned on the
text Jev receives, so a tabular row is its serialized lines. Both select their configuration on
validation log-loss only, are trained without evaluation labels, and are refitted on nested
fractions of the training part for learning curves. No model is recalibrated.

\paragraph{Label budgets and cost.}
Jev and Qwen use no training labels and use validation labels only where a decision or escalation
threshold is selected; XGBoost and BERT use every training and validation label
(Appendix~\ref{app:protocol}, Table~\ref{tab:labels}). All Jev requests cost USD~\JevCostUSD{}
and all cloud training and inference jobs USD~\AwsCostTotalUSD; a fifth dataset was run and then
removed from the study, and its share of both is stated in Appendix~\ref{app:protocol}.
